\subsection{An equivalent inverse without building a matrix}
\label{sec:inverse_without_matrix}
For SM4's particular data transformation $L$, the matrix implementation can be
replaced by three applications of the \emph{forward} map followed by an eight-bit
right rotation:
\[
\boxed{L^{-1}(y)=\operatorname{ROR}_8\bigl(L(L(L(y)))\bigr).}
\]
These are not three decryptions, three fault injections, or three cipher rounds.
They are local evaluations of the same public word transformation.
The existing matrix implementation remains active in \texttt{piret\_sm4.py};
this alternative is provided there as a commented replacement.

\subsubsection*{Why the identity holds for every 32-bit input}
Write $\rho_n$ for a 32-bit left rotation by $n$ bits and $I$ for the identity
map. This notation avoids confusing rotations with the diagram's output words
$R_0,R_1,R_2,R_3$. Rotations commute and compose by adding their amounts:
$\rho_a\rho_b=\rho_{a+b}$, with amounts reduced modulo 32.
The data map is
\[
L=I\xor\rho_2\xor\rho_{10}\xor\rho_{18}\xor\rho_{24}.
\]
When this sum is composed with itself, each cross term occurs twice and cancels
over GF(2). Only the squared terms remain:
\[
\begin{aligned}
L^2
&=I\xor\rho_4\xor\rho_{20}\xor\rho_{36}\xor\rho_{48}\\
&=I\xor\underbrace{\rho_4\xor\rho_4}_{0}\xor\rho_{20}\xor\rho_{16}\\
&=I\xor\rho_{20}\xor\rho_{16}.
\end{aligned}
\]
Apply the same cancellation once more:
\[
\begin{aligned}
L^4&=(L^2)^2=I\xor\rho_{40}\xor\rho_{32}\\
&=\underbrace{I\xor I}_{0}\xor\rho_8=\rho_8.
\end{aligned}
\]
Now let the observed input to the inverse be $y=L(x)$. Three additional
applications give
\[
L^3(y)=L^4(x)=\operatorname{ROL}_8(x).
\]
Rotating right by 8 restores $x$. Because rotations commute with $L$, the same
formula is also a right inverse; it is the unique inverse represented by the
matrix method. This is an algebraic derivation from the definition of $L$,
not a property of arbitrary invertible matrices or the key-schedule map $L'$.

\subsubsection*{Optional Python implementation}
\begin{Verbatim}
def inverse_linear(value):
    # Three applications of the FORWARD data map.
    result = sm4.linear(value)
    result = sm4.linear(result)
    result = sm4.linear(result)

    # A left rotation by 24 is a right rotation by 8 on 32 bits.
    result = sm4.rotl(result, 24)

    return result
\end{Verbatim}
Only one function named \texttt{inverse\_linear} should be active at a time.
Replacing the implementation does not change the candidate equations or the
recovered keys. No inverse S-box is involved in this transformation.

\subsubsection*{Use the same numerical example as the matrix section}
The previous example has $y=\texttt{734E3D3D}$. The following values are
generated and checked against the current cipher functions:
\begin{center}\begin{tabular}{ll}\toprule
Operation & Actual result (hex)\\\midrule
Input $y$ & \texttt{734E3D3D}\\
Forward L application 1 & \texttt{4F04BF00}\\
Forward L application 2 & \texttt{9DA47A90}\\
Forward L application 3 & \texttt{0000004F}\\
Right rotation by 8 & \texttt{4F000000}\\
Existing matrix implementation & \texttt{4F000000}\\
\bottomrule\end{tabular}\end{center}


\begin{center}
\begin{tikzpicture}[node distance=4mm,every node/.style={font=\small},box/.style={draw,rounded corners,align=center,fill=teal!5,minimum height=12mm,text width=25mm}]
\node[box] (a) {\texttt{734E3D3D}\\observed $y$};
\node[box,right=of a] (b) {\texttt{4F04BF00}\\after $L$};
\node[box,right=of b] (c) {\texttt{9DA47A90}\\after $L$};
\node[box,right=of c] (d) {\texttt{0000004F}\\after $L$};
\node[box,right=of d] (e) {\texttt{4F000000}\\right rotate 8};
\draw[-{Stealth}] (a)--(b);\draw[-{Stealth}] (b)--(c);
\draw[-{Stealth}] (c)--(d);\draw[-{Stealth}] (d)--(e);
\end{tikzpicture}
\end{center}
The final rotation is explicitly
\[
\begin{aligned}
\operatorname{ROR}_8(\texttt{0000004F})
&=(\texttt{0000004F}\gg8)
\mathbin{|}\bigl((\texttt{0000004F}\ll24)\mathbin{\&}\texttt{FFFFFFFF}\bigr)\\
&=\texttt{00000000}\mathbin{|}\texttt{4F000000}
=\texttt{4F000000}.
\end{aligned}
\]
Thus both implementations return the \emph{same difference}
$\Delta B=\texttt{4F000000}$, not the individual clean S-box output.
The forward check is $L(\texttt{4F000000})=\texttt{734E3D3D}$.
The builder additionally checks both implementations on all 32 unit vectors;
their linearity makes agreement on this basis sufficient for all 32-bit words.

\note{The matrix version explains how to invert any suitable binary linear map.
The rotation version exploits a special identity of this SM4 map and avoids
building or storing its inverse matrix. Both are correct representations of
the same operation; this guide does not claim a measured performance advantage.}
